\section{方法论：物理限制与图状态机的暴力美学}
\label{sec:methodology}

Grox\_chat 彻底放弃了由单一模型控制所有环节的自由探索范式。我们认为，那些无法独立完成长周期规划的廉价大语言模型，就像是充满创造力但极易走火入魔的“偏执狂”。我们的做法是用代码和图状态机（LangGraph）为它们打造一个具有高度自愈能力的“赛博牢笼”，强行将其纳入正轨。

\subsection{利用自然语言提取意图的双轨 RAG}
让廉价模型按要求输出带有层层嵌套的 API 请求 JSON，无异于一场灾难。这是它们被阻挡在高级 Multi-Agent 门外的核心原因。Grox\_chat 用最粗暴但有效的方式解决了这个问题：“你们不用写代码，说人话就行”。
在每一次智能体发言前，我们会在极短的上下文后加上一句简单的提示：“用一句纯中文说出你现在最想去网上查什么数据。”
底层系统接收到这句简短的自然语言意图后，直接在后台代为执行检索流程：调用 \texttt{all-MiniLM-L6-v2} 将这句中文向量化，接着用 \texttt{sqlite-vec} 在包含了所有被验证过的历史事实数据库中进行 KNN 召回，最后交给 \texttt{bge-reranker-v2-m3} 交叉编码器进行打分过滤。那些得分超过阈值的“硬核事实”，会在该智能体正式发言前，像一张张“小抄”一样被强行塞进它的视野。这种“保姆式”的喂饭服务，让低端模型毫无负担地享受到了 SOTA 级别的信息支撑。

\subsection{异构模型的阶级分化与 SQLite 脑皮层}
在 Grox\_chat 中，模型是分阶级的：
\begin{itemize}
    \item \textbf{高贵的典狱长（Gemini 3.1 Pro/3.0 Flash）}：负责一切不容出错的宏观规划——生成话题大纲、定点撰写 Background Brief、阶段性总结以及判断是否应该强行结题。
    \item \textbf{苦力的服刑者（MiniMax M2.5）}：扮演科学家、空想家、工程师等专家角色，在这个封闭的环形斗兽场里进行廉价且高频的口水战。
\end{itemize}
由于服刑者记忆力差，我们彻底褫夺了它们管理上下文的权力。所有的发言都会被切断、压缩并下沉到 SQLite 数据库中。它们在吵架时，除了刚刚发生的 6 句话和 RAG 贴进来的硬事实，什么都看不见。这保证了哪怕吵到第 100 轮，它们也不会被过去累积的几万字上下文带偏焦点。

\subsection{DAPA 机制：非理性气氛组的降维惩戒}
如果全员都是理性派，当一个模型开始用伪科学胡说八道时，其他廉价模型很容易因为“过度礼貌”而顺着它往下编，这就是最可怕的幻觉级联（Hallucination Cascades）。
为此，Grox\_chat 引入了由独立判定节点组成的**非理性气氛组**，它们执行一套名为动态注意力与事实压力 (Dynamic Attention and Factual Pressure, DAPA) 的机制。
\begin{itemize}
    \item \textbf{狗 (Dot)}：一个只会狂吠的找茬机器。它会在一轮结束后跳出来，死死咬住本轮最像胡编乱造的那个人。
    \item \textbf{猫 (Cat)}：只凭喜好奖励本轮最有趣的灵魂。
    \item \textbf{裁决者 (Tron: Fight for user)}：一个严格执行阿西莫夫式“论坛三定律”的冰冷裁判，专门审查是否出现了反人类、极度偏置或严重的学术造假。
\end{itemize}
一旦有人被狗咬或被 Tron 警告，LangGraph 的状态图就会发生物理形变：系统会在所有专家发言完毕后，**强行凭空插队**出一个额外的“惩罚回合（Penalty Turn）”。在这个回合里，那个被锁定的倒霉蛋将面临来自系统底层的无情斥责，并被迫利用 RAG 数据对刚刚说出口的胡话进行公开处刑和逻辑推翻。通过这种简单粗暴的图结构追加，系统获得了惊人的“自愈”能力。